\textbf{Phần 1. Trắc nghiệm lựa chọn}
\setcounterpageref{ex}{0}

\begin{ex} %[5D6N31]
Kết quả của phép tính $7{,}8 + 11{,}6$ là
\choice
{\True $19{,}4$}
{$18{,}4$}
{$19{,}14$}
{$18{,}14$}
\loigiai{
Ta có: $7{,}8 + 11{,}6 = 19{,}4$.\\
\textbf{Đáp án của câu hỏi là A}
}
\end{ex}

\begin{ex} %[5D6N31]
Kết quả của phép tính $27 - 1{,}2$ là
\choice
{$15{,}8$}
{$26{,}8$}
{\True $25{,}8$}
{$25{,}2$}
\loigiai{
Ta có: $27{,}0 - 1{,}2 = 25{,}8$.\\
\textbf{Đáp án của câu hỏi là C}
}
\end{ex}

\begin{ex} %[5D6H31]
Tính thuận tiện biểu thức $9{,}2 + 17{,}56 + 0{,}8$ ta được kết quả là
\choice
{$26{,}56$}
{$28{,}56$}
{$37{,}56$}
{\True $27{,}56$}
\loigiai{
Ta có: $(9{,}2 + 0{,}8) + 17{,}56 = 10 + 17{,}56 = 27{,}56$.\\
\textbf{Đáp án của câu hỏi là D}
}
\end{ex}

\begin{ex} %[5D6H31]
Kết quả của phép tính $9{,}63 - 5{,}275$ là
\choice
{\True $4{,}355$}
{$4{,}365$}
{$4{,}455$}
{$4{,}445$}
\loigiai{
Ta có: $9{,}630 - 5{,}275 = 4{,}355$.\\
\textbf{Đáp án của câu hỏi là A}
}
\end{ex}

\begin{ex} %[5D6H31]
Cho biết $29 \times 37 = 1\,073$. Giá trị của phép tính $0{,}29 \times 3{,}7$ là
\choice
{$107{,}3$}
{$10{,}73$}
{\True $1{,}073$}
{$0{,}1073$}
\loigiai{
Hai thừa số có tất cả $2 + 1 = 3$ chữ số ở phần thập phân nên kết quả là $1{,}073$.\\
\textbf{Đáp án của câu hỏi là C}
}
\end{ex}

\begin{ex} %[5D6H31]
Kết quả của phép tính $1{,}2 \times 3{,}5$ là
\choice
{$0{,}42$}
{$42$}
{$4{,}25$}
{\True $4{,}2$}
\loigiai{
Ta có: $1{,}2 \times 3{,}5 = 4{,}2$.\\
\textbf{Đáp án của câu hỏi là D}
}
\end{ex}

\begin{ex} %[5D6H31]
Ba bạn gom giấy vụn: Mai gom $2{,}5\text{ kg}$, Nam gom $3{,}25\text{ kg}$, Việt gom $2{,}4\text{ kg}$. Bạn gom được nhiều giấy vụn nhất hơn bạn gom được ít giấy vụn nhất là
\choice
{$0{,}75\text{ kg}$}
{\True $0{,}85\text{ kg}$}
{$0{,}1\text{ kg}$}
{$1{,}15\text{ kg}$}
\loigiai{
Nhiều nhất là Nam ($3{,}25\text{ kg}$), ít nhất là Việt ($2{,}4\text{ kg}$).\\
Hiệu: $3{,}25 - 2{,}4 = 0{,}85\text{ (kg)}$.\\
\textbf{Đáp án của câu hỏi là B}
}
\end{ex}

\begin{ex} %[5D6V31]
Rô-bốt cao $0{,}9\text{ m}$. Bạn Mi cao hơn Rô-bốt $0{,}35\text{ m}$ và thấp hơn Mai $0{,}31\text{ m}$. Chiều cao của bạn Mai là
\choice
{$1{,}25\text{ m}$}
{$1{,}55\text{ m}$}
{\True $1{,}56\text{ m}$}
{$1{,}58\text{ m}$}
\loigiai{
Mi cao: $0{,}9 + 0{,}35 = 1{,}25\text{ (m)}$.\\
Mai cao: $1{,}25 + 0{,}31 = 1{,}56\text{ (m)}$.\\
\textbf{Đáp án của câu hỏi là C}
}
\end{ex}

\begin{ex} %[5D6V31]
Trên hai đĩa cân đĩa thăng bằng: Đĩa trái có con Thỏ và hai quả cân $2\text{ kg}, 500\text{ g}$, đĩa phải có con Cáo. Ở một cân khác, quả cân $5\text{ kg}$ thăng bằng với con Cáo và quả cân $200\text{ g}$. Cân nặng của con Thỏ là
\choice
{$4{,}8\text{ kg}$}
{\True $2{,}3\text{ kg}$}
{$2{,}5\text{ kg}$}
{$2{,}8\text{ kg}$}
\loigiai{
Cân thứ hai cho biết: Cáo $+ 0{,}2\text{ kg} = 5\text{ kg} \Rightarrow$ Cáo nặng $5 - 0{,}2 = 4{,}8\text{ kg}$.\\
Cân thứ nhất cho biết: Thỏ $+ 2{,}5\text{ kg} = \text{Cáo} \Rightarrow$ Thỏ nặng $4{,}8 - 2{,}5 = 2{,}3\text{ kg}$.\\
\textbf{Đáp án của câu hỏi là B}
}
\end{ex}

\begin{ex} %[5D6V31]
Sử dụng các tấm thẻ số $6$; $3$; $5$ và một dấu phẩy $,$, hãy lập số thập phân lớn nhất và số thập phân bé nhất (mỗi thẻ dùng đúng một lần). Hiệu của hai số thập phân đó là
\choice
{$61{,}74$}
{$31{,}64$}
{$59{,}74$}
{\True $61{,}74$}
\loigiai{
Số thập phân lớn nhất có thể lập là $65{,}3$.\\
Số thập phân bé nhất có thể lập là $3{,}56$.\\
Hiệu hai số là: $65{,}30 - 3{,}56 = 61{,}74$.\\
\textbf{Đáp án của câu hỏi là D}
}
\end{ex}

\begin{ex} %[5D6V31]
Một mảnh đất hình chữ nhật có chiều dài $24{,}5\text{ m}$ và chiều rộng kém chiều dài $6{,}5\text{ m}$. Chu vi của mảnh đất đó là
\choice
{\True $85\text{ m}$}
{$18\text{ m}$}
{$42{,}5\text{ m}$}
{$441\text{ m}$}
\loigiai{
Chiều rộng mảnh đất là: $24{,}5 - 6{,}5 = 18\text{ (m)}$.\\
Chu vi mảnh đất là: $(24{,}5 + 18) \times 2 = 42{,}5 \times 2 = 85\text{ (m)}$.\\
\textbf{Đáp án của câu hỏi là A}
}
\end{ex}

\begin{ex} %[5D6V31]
Một cửa hàng ngày thứ nhất bán được $125{,}5\text{ kg}$ gạo, ngày thứ hai bán được nhiều hơn ngày thứ nhất $34{,}8\text{ kg}$ gạo. Cả hai ngày cửa hàng bán được tất cả là
\choice
{$160{,}3\text{ kg}$}
{\True $285{,}8\text{ kg}$}
{$286{,}8\text{ kg}$}
{$275{,}8\text{ kg}$}
\loigiai{
Ngày thứ hai bán được: $125{,}5 + 34{,}8 = 160{,}3\text{ (kg)}$.\\
Cả hai ngày bán được: $125{,}5 + 160{,}3 = 285{,}8\text{ (kg)}$.\\
\textbf{Đáp án của câu hỏi là B}
}
\end{ex}

\textbf{Phần 2. Trắc nghiệm đúng sai}
\setcounterpageref{ex}{0}

\begin{ex} %[5D6H31]
Xét tính đúng sai của các phép tính sau:
\choiceTF
{\True Phép cộng $92{,}15 + 7{,}99 + 0{,}01 = 100{,}15$.}
{Phép trừ $19{,}3 - 11 = 18{,}2$.}
{\True Phép nhân $3{,}1 \times 5{,}7 = 17{,}67$.}
{\True Phép nhân $2{,}9 \times 3{,}7 = 10{,}73$.}
\loigiai{
\begin{itemize}
    \item $92{,}15 + (7{,}99 + 0{,}01) = 92{,}15 + 8 = 100{,}15$.
    \item $19{,}3 - 11 = 8{,}3 \neq 18{,}2$.
    \item $3{,}1 \times 5{,}7 = 17{,}67$.
    \item $2{,}9 \times 3{,}7 = 10{,}73$.
\end{itemize}
\begin{itemize}
    \item a) \textbf{Đúng}.
    \item b) \textbf{Sai}: Kết quả là $8{,}3$.
    \item c) \textbf{Đúng}.
    \item d) \textbf{Đúng}.
\end{itemize}
}
\end{ex}

\begin{ex} %[5D6V31]
Ba bạn Mai, Nam, Việt thi gom giấy vụn với kết quả: Mai gom $2{,}5\text{ kg}$, Nam gom $3{,}25\text{ kg}$, Việt gom $2{,}4\text{ kg}$. Xét tính đúng sai của các khẳng định sau:
\choiceTF
{\True Nam là bạn gom được nhiều giấy vụn nhất.}
{Tổng số giấy vụn cả $3$ bạn gom được là $7{,}15\text{ kg}$.}
{\True Nam gom nhiều hơn Mai $0{,}75\text{ kg}$ giấy vụn.}
{\True Tổng khối lượng giấy của Mai và Việt gom được nhiều hơn Nam.}
\loigiai{
\begin{itemize}
    \item Nam gom $3{,}25\text{ kg}$ là nhiều nhất.
    \item Tổng số giấy: $2{,}5 + 3{,}25 + 2{,}4 = 8{,}15\text{ kg} \neq 7{,}15\text{ kg}$.
    \item Nam nhiều hơn Mai: $3{,}25 - 2{,}5 = 0{,}75\text{ kg}$.
    \item Mai + Việt $= 2{,}5 + 2{,}4 = 4{,}9\text{ kg} > 3{,}25\text{ kg}$ (Nam).
\end{itemize}
\begin{itemize}
    \item a) \textbf{Đúng}.
    \item b) \textbf{Sai}: Tổng là $8{,}15\text{ kg}$.
    \item c) \textbf{Đúng}.
    \item d) \textbf{Đúng}.
\end{itemize}
}
\end{ex}

\textbf{Phần 3. Trả lời ngắn}
\setcounterpageref{ex}{0}

\begin{ex} %[5D6H31]
Tính giá trị biểu thức: $9{,}2 + 17{,}56 + 0{,}8$. (Nhập kết quả dưới dạng số thập phân).
\shortans{27,56}
\loigiai{
\textbf{Đáp số:} $27{,}56$\\
$(9{,}2 + 0{,}8) + 17{,}56 = 10 + 17{,}56 = 27{,}56$.
}
\end{ex}

\begin{ex} %[5D6H31]
Tính kết quả phép tính: $9{,}63 - 5{,}275$. (Nhập kết quả dưới dạng số thập phân).
\shortans{4,355}
\loigiai{
\textbf{Đáp số:} $4{,}355$\\
$9{,}63 - 5{,}275 = 4{,}355$.
}
\end{ex}

\begin{ex} %[5D6H31]
Tính kết quả phép tính: $1{,}2 \times 3{,}5$. (Nhập kết quả dưới dạng số thập phân).
\shortans{4,2}
\loigiai{
\textbf{Đáp số:} $4{,}2$\\
$1{,}2 \times 3{,}5 = 4{,}2$.
}
\end{ex}

\begin{ex} %[5D6V31]
Rô-bốt cao $0{,}9\text{ m}$, Mi cao hơn Rô-bốt $0{,}35\text{ m}$, Mai cao hơn Mi $0{,}31\text{ m}$. Hỏi bạn Mai cao bao nhiêu mét? (Nhập kết quả dưới dạng số thập phân).
\shortans{1,56}
\loigiai{
\textbf{Đáp số:} $1{,}56$\\
$0{,}9 + 0{,}35 + 0{,}31 = 1{,}56\text{ (m)}$.
}
\end{ex}

\begin{ex} %[5D6V31]
Biết quả cân $5\text{ kg}$ thăng bằng với con Cáo và quả cân $0{,}2\text{ kg}$; con Cáo lại thăng bằng với con Thỏ và quả cân $2{,}5\text{ kg}$. Tính cân nặng của con Thỏ theo ki-lô-gam. (Nhập kết quả dưới dạng số thập phân).
\shortans{2,3}
\loigiai{
\textbf{Đáp số:} $2{,}3$\\
Cáo: $5 - 0{,}2 = 4{,}8\text{ kg}$. Thỏ: $4{,}8 - 2{,}5 = 2{,}3\text{ kg}$.
}
\end{ex}

\begin{ex} %[5D6V31]
Cho các thẻ số $6, 3, 5$ và dấu phẩy. Lập số lớn nhất và số bé nhất rồi tính hiệu của hai số đó. (Nhập kết quả dưới dạng số thập phân).
\shortans{61,74}
\loigiai{
\textbf{Đáp số:} $61{,}74$\\
$65{,}3 - 3{,}56 = 61{,}74$.
}
\end{ex}
